% =============================================================================
% CAPÍTULO 7 — Conclusiones Generales y Lecciones Aprendidas
% =============================================================================
\chapter{Conclusiones Generales y Lecciones Aprendidas}
\label{cap:conclusiones}

Este capítulo cierra la memoria consolidando los logros del desarrollo de \textbf{EthosHub} a lo
largo de los seis incrementos (Sprint~0 a Sprint~5), las lecciones aprendidas y las líneas de
trabajo futuro, conforme al cierre de versión documentado en el Capítulo~\ref{cap:sprint5}.\par

\section{Logros del Producto}
\label{sec:concl_logros}
El proyecto entregó una plataforma funcional de portafolios profesionales y descubrimiento de
talento, con un despliegue monolítico verificable. Los logros más relevantes son:

\begin{itemize}[noitemsep]
    \item \textbf{Seguridad delegada y robusta:} identidad en Supabase Auth, validación de JWT
    como \textit{Resource Server}, OTP para operaciones sensibles y mitigaciones OWASP
    (Capítulos~\ref{cap:sprint1} y~\ref{cap:sprint4}).
    \item \textbf{Dominio completo:} perfil, trayectoria, habilidades, proyectos con evidencias y
    conexiones externas, todo aislado por RLS (Capítulos~\ref{cap:sprint2} y~\ref{cap:sprint3}).
    \item \textbf{Funcionalidades en tiempo real:} chat con presencia e idempotencia, pipeline
    Kanban e \textit{insights} de IA (Capítulo~\ref{cap:sprint3}).
    \item \textbf{Calidad y despliegue:} pruebas E2E, accesibilidad y empaquetado monolítico
    productivo (Capítulos~\ref{cap:sprint4} y~\ref{cap:sprint5}).
\end{itemize}

\section{Lecciones Aprendidas}
\label{sec:concl_lecciones}
\begin{enumerate}[noitemsep]
    \item \textbf{Seguridad por diseño:} canalizar las escrituras sensibles por RPCs
    \comando{SECURITY DEFINER} con \comando{auth.uid()} evitó fugas de datos entre usuarios.
    \item \textbf{Idempotencia temprana:} introducir \comando{client\_message\_id} desde el
    diseño del chat eliminó los duplicados bajo reconexión.
    \item \textbf{Documentation as Code:} versionar la documentación junto al código mantuvo la
    información sincronizada con el producto, cumpliendo ISO 26515.
    \item \textbf{Deuda técnica acotada:} planificar su reducción en el Sprint de QA evitó que se
    acumulara hasta el lanzamiento.
\end{enumerate}

\section{Trabajo Futuro}
\label{sec:concl_futuro}
\begin{itemize}[noitemsep]
    \item Ampliar las integraciones con \textit{providers} externos (importación automática).
    \item Profundizar los \textit{insights} de IA con más modos de análisis.
    \item Incorporar observabilidad avanzada (trazas distribuidas) y \textit{dashboards} de SLO.
    \item Externalizar y rotar de forma definitiva todos los secretos del historial de git.
\end{itemize}

\section{Cumplimiento Normativo}
\label{sec:concl_iso}
La memoria evidencia el cumplimiento de \textbf{ISO/IEC/IEEE 26515:2024} (información en entornos
ágiles) y \textbf{ISO/IEC/IEEE 15289:2024} (productos de información), según la matriz de
cumplimiento del Anexo~\ref{cap:anexo_iso}. Cada actividad y producto de información exigido tiene
una evidencia trazable en este documento o en sus anexos, versionada junto al código fuente.

\section{Indicadores Globales del Proyecto}
\label{sec:concl_indicadores}
La tabla~\ref{tab:concl_indicadores} consolida los indicadores del desarrollo completo de
EthosHub, desde la incepción hasta el cierre de la versión 2.0.0.

\begin{table}[!ht]
    \centering
    \renewcommand{\arraystretch}{1.3}
    \fontsize{9}{10.8}\selectfont\sffamily
    \begin{tabularx}{\textwidth}{|X|l|}
        \hline
        \rowcolor{colorPrimario}
        \textcolor{white}{\textbf{Indicador}} & \textcolor{white}{\textbf{Valor}} \\ \hline
        Incrementos entregados (Sprint 0--5) & 6 \\ \hline
        \rowcolor{grisTecnico}
        Duración del desarrollo (Sprints 1--5) & 10 semanas (30 mar--6 jun 2026) \\ \hline
        Épicas completadas & 8 \\ \hline
        \rowcolor{grisTecnico}
        Migraciones de base de datos & V1--V15 \\ \hline
        Dominios funcionales del backend & 16 \\ \hline
        \rowcolor{grisTecnico}
        Versión final del producto & \comando{ethoshub:2.0.0} \\ \hline
        Estándares cumplidos & ISO/IEC/IEEE 26515 y 15289 \\ \hline
    \end{tabularx}
    \caption{Indicadores globales del proyecto EthosHub.}
    \label{tab:concl_indicadores}
\end{table}

Con el cierre del Sprint~5, EthosHub alcanza un producto desplegable, documentado y trazable, que
satisface los compromisos de entrega del Plan General y los criterios de aceptación fijados en la
incepción (Sección~\ref{ssec:s0_dod}).

\clearpage
